% Source-aligned English for printed p.133.
\SourcePage{133}

\LayerTarget{tohoku:example:1.7:h:p2}
\noindent corresponding diagrams, for a given category $C$ not necessarily
additive, are known as \textit{inductive systems} on $I$ with values in $C$.
Replacing $I$ by the opposite preordered set, or $C$ by $C^\circ$, gives the
\textit{projective systems} on $I$ with values in $C$. An important case is
that in which $I$ is the set of open subsets of a topological space $X$,
ordered by the relation $\supset$; this gives the notion of a
\textit{presheaf} on $X$ with values in the category $C$.

\LayerTarget{tohoku:section:1.8}
\noindent\textbf{1.8. Inductive and projective limits.} We shall discuss only
the former, since the notion of a projective limit is dual to that of an
inductive limit. Let $C$ be a category, let $I$ be a preordered set, and let
$\mathbf A=(A_i,u_{ij})$ be an inductive system on $I$ with values in $C$
($u_{ij}$ is a morphism $A_j\longrightarrow A_i$ defined for $i\geq j$).

\LayerTarget{tohoku:def:inductive-limit}
A (generalized) \textit{inductive limit} of $\mathbf A$ is a system consisting
of an object $\overline A\in C$ and a family $(u_i)$ of morphisms
$u_i:A_i\longrightarrow\overline A$ satisfying the following conditions:
a) if $i\leq j$, then $u_i=u_j u_{ji}$; b) for every $B\in C$ and every family
$(v_i)$ of morphisms $v_i:A_i\longrightarrow B$ such that
$v_i=v_j u_{ji}$ whenever $i\leq j$, there is a unique morphism $v$ from
$\overline A$ to $B$ such that $v_i=v u_i$ for every $i\in I$.

\LayerTarget{tohoku:def:inductive-limit:functoriality}
Suppose $(\overline A,(u_i))$ is an inductive limit of
$\mathbf A=(A_i,u_{ij})$, $(\overline B,(v_i))$ is an inductive limit of a
second inductive system $\mathbf B=(B_i,v_{ij})$, and $w=(w_i)$ is a morphism
from $\mathbf A$ to $\mathbf B$. Then there is a unique morphism
$\overline w$ from $\overline A$ to $\overline B$ such that
$\overline w u_i=v_iw_i$ for every $i\in I$. In particular, any two inductive
limits of the same inductive system are canonically isomorphic, in the evident
sense. It is therefore natural to choose, for every inductive system admitting
an inductive limit, one such limit (for example by means of Hilbert's
$\tau$-symbol), denoted
$\mathop{\rm lim}\limits_{\longrightarrow}\mathbf A$ or
$\mathop{\rm lim}\limits_{\longrightarrow,\,i\in I}A_i$, and
called \textit{the} inductive limit of the given inductive system.

\LayerTarget{tohoku:def:inductive-limit:functor}
If $I$ and $C$ are such that
$\mathop{\rm lim}\limits_{\longrightarrow}\mathbf A$ exists for
\textit{every} inductive system $\mathbf A$ on $I$ with values in $C$, it
follows that $\mathop{\rm lim}\limits_{\longrightarrow}\mathbf A$ is a
\textit{covariant functor} from the category of inductive systems from $I$ to
$C$, with values in $C$.

\LayerTarget{tohoku:prop:1.8}
\noindent\textsc{Proposition 1.8.} --- \textit{Let $C$ be an abelian category
satisfying axiom AB 3), the existence of arbitrary direct sums, and let $I$ be
an increasing filtered preordered set. Then for every inductive system
$\mathbf A$ on $I$ with values in $C$, the inductive limit
$\mathop{\rm lim}\limits_{\longrightarrow}\mathbf A$ exists, and as a functor
of $\mathbf A$ it is additive and right exact. If $C$ satisfies axiom AB 5)
(cf. 1.5), this functor is even exact, and the kernel of the canonical
morphism
$u_i:A_i\longrightarrow\mathop{\rm lim}\limits_{\longrightarrow}\mathbf A$
is then the supremum of the kernels of the morphisms
$u_{ji}:A_i\longrightarrow A_j$ for $j\geq i$. In particular, if the $u_{ji}$
are injective, then so is $u_i$.}

\LayerTarget{tohoku:prop:1.8:proof:p1}
To construct an inductive limit of $(A_i,u_{ij})$, set

\[
  S=\bigoplus_{i\in I}A_i,
\]

\noindent and, for every pair $(i,j)$ with $i\leq j$, let $w_{ij}$ be the
morphism from $A_i$ to $S$ deduced from the morphism
$A_i\longrightarrow A_i+A_j$ whose components are $I_{A_i}$ and $-u_{ji}$.
Let $N_{ij}$

\TohokuContinuation{134}
